\section{Related Work}

\subsection{Selective Classification and Group Fairness}
Geifman and El-Yaniv showed how to add a reject option to an already-trained deep network by choosing a confidence threshold so that the error on accepted cases stays below a target level with high probability \cite{geifman2017}. Such thresholds are set on average performance. Jones et al. found that, while abstention improves average accuracy, it can magnify existing accuracy gaps between groups and, in some groups, lower accuracy as abstention increases \cite{jones2021}. They linked this to the distribution of model confidence and proposed training models to perform equally across groups at full coverage. Their analysis treats the full-coverage gap as given and does not ask where it comes from, which is the question we address. Closer to deployment, a recent preprint proposes a fairness-aware framework that routes medical images to a model, a human expert, or both, and accounts for subgroup-dependent reliability without using sensitive attributes at test time \cite{zhang2026}. Such systems still rely on some signal of when the model is unreliable, and our results bear on when confidence fails to provide it.

\subsection{Annotation Bias and Subgroup Separability in Chest Radiography}
Labels in large public chest X-ray datasets are mostly extracted from radiology reports. Seyyed-Kalantari et al. trained classifiers on three such datasets and found that female, young, Black, Hispanic and Medicaid-insured patients were more often wrongly labelled healthy, and argued that labels drawn from clinical records can carry existing biases in care \cite{seyyedkalantari2021}. Santomartino et al. evaluated four report-labelling tools and found high pooled accuracy but error rates that varied with patient age, with most tools performing worst in patients over 80 \cite{santomartino2024}. Jones, Roschewitz and Glocker related subgroup separability, meaning how easily a model can tell subgroups apart, to the disparities and performance loss that follow training on systematically biased labels such as underdiagnosis \cite{jones2023}. We use separability as one of our four factors and test whether it also governs how label bias reaches abstention.

\subsection{Domain Shift and Hidden Failures}
Zech et al. found that pneumonia models often performed worse on radiographs from external hospital systems than on new data from the training hospitals, and that the models could identify which hospital and department an image came from \cite{zech2018}. Oakden-Rayner et al. showed that medical imaging models can look accurate overall while failing on unlabelled subsets of patients, and argued that such hidden stratification should be checked before deployment \cite{oakdenrayner2020}. Rafferty and Rajan, in a preprint, reviewed label quality, bias and domain shift in public chest radiography datasets and reported substantial loss of performance in cross-dataset testing \cite{rafferty2025}. These studies measure accuracy under shift; we ask how shift changes the confidence that abstention relies on, and therefore whether abstention signals the harm.
